\section{Polynomial time and the output string}\label{sec:algorithm}

\begin{proof}[Proof of \Cref{thm:main}]
Parse and preprocess the explicit input in time polynomial in its encoded size $N$. This handles the empty instance and preserves the optimum otherwise. In the nonempty case, construct the counts and base graph of \Cref{sec:counts}. By \Cref{lem:counts,lem:base}, its cost is $W\le\OPT(\mathcal S)$ and it contains every required vertex. Apply \Cref{thm:connection}. The resulting graph is balanced, its used support is connected, it includes $\eps$, and its total up cost is at most $2W$.

A finite balanced directed multigraph with connected used support has a closed Euler tour from any used vertex. One may obtain it by following unused edges until a closed walk is formed, and splicing further closed walks at visited vertices until all edges have been used. Balance prevents an unfinished walk from stopping away from its starting point, and connectivity ensures that any unused edges can eventually be reached for a splice.

Start such a tour at $\eps$ and write down the appended letter at each up step. At every step, the current vertex word is a suffix of the string written so far: an up step appends a letter to both, and a down step removes the first letter of the vertex. Every required string is visited and therefore occurs as a substring of the output. Its length is the number of up steps, at most $2W\le2\OPT(\mathcal S)$.

We finish by verifying that all preceding constructions are deterministic polynomial-time operations on explicit finite data. Order encoded symbols, input strings, substring vertices, layers, records, and edge copies by their explicit codes and stable indices. At every finite choice in the construction, including walk decompositions, eligible records, tied extrema, and Euler splices, take the first admissible item in this order. Comparisons and output of variable-width symbol codes are charged their bit cost; no unit-cost operation on an unbounded label or random choice is used.

\paragraph{Counts and graph.}
There are $O(L^2)$ words in $V$, each of length at most $L\le N$; store their explicit labels or interned identifiers in a dictionary. For each $v$, enumerate $1\le p<\len v$ and $k\ge1$ with $kp<\len v$. The period-$p$ template is specified by $v[0:p)$, and testing whether it matches $v$ takes polynomial time. Thus all rules \eqref{eq:period-rule} form a polynomial list. Each blocking functional is evaluated by enumerating $r\in V$ and the possible placements in \eqref{eq:block-interior}--\eqref{eq:block-ends}. The descending count recursion uses only such evaluations and integer maxima. By \Cref{lem:counts}, all counts are at most $L$, so their bit lengths are polynomial as well.

By \Cref{lem:base}, the explicit base graph has exactly $2W\le2L$ edges with multiplicity; adjacency lists retain individual edge-copy IDs. Its closed-walk decomposition can be found by finite edge traversals. Primitive periods, alignments of texts, and sorted passage lists can be found by direct string comparisons and sorting. The total number of primitive layer phases is the sum of their periods, namely $W$. Store each layer by its primitive period word, one period of exit-end values $z(x)$, and a stable ID; \eqref{eq:layer-properties} recovers other coordinates without enumerating an infinite text.

\paragraph{Records and connections.}
The first and last ends of a layer, its minimum vertex length, and the distinguished rotation of its text can be computed from one period. Equality or lexicographic comparison of two periodic texts requires at most the sum of their periods to find a difference, by \Cref{lem:agreement}. Normalize the start of each considered periodic record to one turn, translating its aligned child coordinates with it as in \Cref{lem:host}. Each actual window labels a substring of length at most $L$, so the resulting coordinates, lengths, and end differences have polynomial bit length. A record is stored by layer ID, phase, normalized start and end, and its alignment.

The search for a record in \Cref{sec:processing} is also finite. A period-$p$ layer has $p$ up steps and $p$ down steps, hence $2p$ window visits per turn, counting the initial vertex but not its repeated final copy. There are therefore $2W$ such visits over all layers. For the specified $w$ and template $A$, enumerate these actual windows and all placements of $w$ in them. Align each placement with $[t,H-K)$, reject a candidate whose complete aligned text equals $A$, and retain one whose entire window agrees with $A$. There are at most $2WL$ candidates. The count rule and \Cref{lem:choose-k,lem:matching-window} guarantee a covering window on a different aligned text, so this search succeeds.

The host operations, collective links, and internal joins require only ordered lists of windows and the paths of \Cref{lem:ordered-windows}; materialize these as explicit edge lists. Finding the intersections used in \Cref{lem:collective} is a scan along at most $K=kp$ units of start coordinate, with $K$ no larger than the group's budget and hence at most $W$. Group processing and the final outgoing-block functional graph likewise use polynomially many layers and blocks.

\paragraph{Materializing the answer.}
Every added walk is closed, so its number of down steps equals its number of up steps. The total added up cost is at most $W$, and hence the final graph has at most $4W\le4L$ edges with multiplicity. Its explicit construction and Euler traversal are polynomial. This proves both the approximation guarantee and the running-time assertion.
\end{proof}
